\subsection{The whole combined Kummer class is excluded by an actual pro-two metabelian quotient}

\begin{statement}{A.157}{Proved}
{A single combined mod-eight class with no ordinary dyadic Selmer image}
Retain \(a\), \(D^\circ\), \(\overline D\), and \(E_b/K\)
from A.154.  Write \(E_U=\pi_1(U/\Q)\),
\(\Pi=\pi_1(U_{\overline\Q})\), and let
\(H\subset\Pi\) be the geometric subgroup of the actual cover
\[
 D^\circ:\ W^8=\frac{t^2(1-t)}{64a^2}\longrightarrow U.
\tag{A.157.1}
\]
Let \(\mathcal I\) be the closed normal subgroup of \(H\)
generated by its four boundary inertia groups, and define
\[
 N=\ker\{H\longrightarrow(H/\mathcal I)^{\rm ab,(2)}\}.
\tag{A.157.2}
\]
Here \({\rm ab,(2)}\) denotes the maximal abelian pro-two
quotient.  The subgroup \(N\) is normal in the full arithmetic
group \(E_U\).  Its geometric quotient has the explicit form
\[
 \begin{gathered}
 Q_\infty=\Pi/N\simeq
       R\rtimes C_8,\qquad
 R=\Z_2[V]/(V^6+V^4+V^2+1),\\
 v\cdot r=Vr,\qquad
 x\longmapsto(1,v^2),\quad y\longmapsto(0,v),
 \end{gathered}
\tag{A.157.3}
\]
where \(x,y\) are meridians about zero and one.
In particular \(R\) is free of rank six over \(\Z_2\),
and \(Q_\infty\) is pro-two and metabelian.

There is no continuous section of \(E_U/N\to G_\Q\)
satisfying all three following conditions: its combined
Kummer class is
\[
          2\kappa_8(t)+\kappa_8(1-t)=[64a^2]
                            \quad\text{in }H^1(G_\Q,\mu_8);
\tag{A.157.4}
\]
at every place its image is geometrically conjugate to
the image of an actual ordinary or cuspidal local
section; and at two it is the image of an ordinary
point section.  The real local condition is included.
No ordinary centre at any prime is prescribed.

Let \(\Pi_2\) be the maximal pro-two quotient of \(\Pi\),
and let \(M_2\) be the inverse image of \(\Pi_2''\).
The actual arithmetic maps
\[
 E_U/\Pi''\longrightarrow E_U/M_2
                         \longrightarrow E_U/N
\tag{A.157.5}
\]
give the same exclusion for locally point-or-cusp-matching
sections of the full metabelian quotient and of the
full pro-two metabelian quotient.  In particular every
ordinary completed unit-equation pair \((A,B)\) satisfying
\([A^2B]_8=[64a^2]\) is excluded from a locally matching
full metabelian Selmer section, and from a full Selmer
section of the tripod.  This applies to all the pairs
of A.145 without its nine prescribed local centres.

The rational points on the actual proper cover are also
determined:
\[
        \overline D(K)=\overline D(\Q)
             =\{P_0,P_1,P_{1/2},P_\infty\}.
\tag{A.157.6}
\]
\end{statement}

\begin{proof}
\emph{The cover and the arithmetic normal subgroup.}
The exact curve maps, including their inverse, are
\[
 \begin{gathered}
 aW^4=T(T-1)(T-\tfrac12),\qquad
 t=4T(1-T),\qquad
 T=\tfrac12-\frac{4aW^4}{t}.
 \end{gathered}
\tag{A.157.7}
\]
Both compositions and the eighth-power equation follow
from \(T(T-1)=-t/4\) and
\((T-1/2)^2=(1-t)/4\), as in A.139.
The geometric Kummer character has values two and one
on \(x,y\), so its image is all of \(\Z/8\), and
\[
 H=\ker\{\Pi\xrightarrow{\,2\kappa_8(t)+\kappa_8(1-t)\,}
                                      \Z/8(1)\}.
\tag{A.157.8}
\]
For an arithmetic element \(e\), conjugation multiplies
this geometric character by \(\chi_8(e)\).
Its kernel is consequently normal in \(E_U\).

The open immersion induces the actual quotient
\[
 H/\mathcal I=\pi_1(\overline D_{\overline\Q}).
\tag{A.157.9}
\]
Indeed a finite cover of \(D^\circ\) on which the
four inertia groups act trivially extends, by normalization,
to a finite cover unramified at all four missing smooth
points.  It is finite etale on \(\overline D\).
Conversely a finite etale cover of the proper curve has
trivial boundary inertia.  This gives (A.157.9) through
the finite-cover correspondence.

Conjugation by \(E_U\) preserves \(\mathcal I\).
One can see the boundary permutation literally on the
geometric deck transformations.  For \(\zeta\in\mu_8\)
they are
\[
 (T,W)\longmapsto
 \begin{cases}(T,\zeta W),&\zeta^4=1,\\
               (1-T,\zeta W),&\zeta^4=-1.
 \end{cases}
\tag{A.157.10}
\]
They extend to the proper normalization and permute
its four boundary points.  Field conjugation has the
same property, because the cover and its boundary are
defined over \(\Q\).  Inner changes preserve the
normal closure of these inertia groups.  Thus the
characteristic abelian pro-two quotient of (A.157.9)
defines an arithmetic normal kernel \(N\), proving
the normality assertion in (A.157.2).

\emph{The integral module and the geometric quotient.}
Put \(u=xy^{-2}\), \(v=y\).  This is an actual Nielsen
basis of the free geometric group, with inverse
\(x=uv^2,y=v\).  Its image modulo \(H\) is
\(u\mapsto0,v\mapsto1\).  The Schreier basis of \(H\) is
\[
 u_j=v^juv^{-j}\quad(0\leq j\leq7),\qquad z=v^8.
\tag{A.157.11}
\]
Its abelianization is free on these nine generators.
The meridians above zero are \(x^4\) and
\(yx^4y^{-1}\); the meridian above one is \(y^8\),
and that above infinity is \((xy)^{-8}\).
These words follow from the inertia orbits of the
degree-eight Kummer cover: the respective orbit
lengths are \(4,4,8,8\).
In the basis (A.157.11) their exact abelian vectors are
\[
 \begin{split}
 [y^8]&=z,\\
 [x^4]&=u_0+u_2+u_4+u_6+z,\\
 [yx^4y^{-1}]&=u_1+u_3+u_5+u_7+z,\\
 [(xy)^{-8}]&=-\sum_{j=0}^7u_j-3z.
 \end{split}
\tag{A.157.12}
\]
For example \(xy=uv^3\), and its eighth power visits
each residue class modulo eight once and has terminal
factor \(v^{24}=z^3\).  The four displayed vectors
sum to zero with the indicated orientations.
Their span is exactly the primitive rank-three module
generated by \(z,\sum_{j\ {
m even}}u_j,
\sum_{j\ {
m odd}}u_j\).

After passage to the abelian pro-two quotient, identify
\(u_j\) with \(V^j\).  Killing these vectors gives
\[
 H/N=\Z_2[V]/(1+V^2+V^4+V^6)=R.
\tag{A.157.13}
\]
The relation \(V^8-1=0\) follows by multiplying the
displayed monic polynomial by \(V^2-1\).
The lift \(v\) has exact order eight, since \(v^8=z\)
is killed and its image in \(\Pi/H\) has that order.
Its conjugation on \(R\) is multiplication by \(V\).
Thus it supplies the actual splitting and maps in
(A.157.3); their inverse sends the module generator
to \(xy^{-2}\) and the cyclic generator to \(y\).

Since \(\Pi/H\) is cyclic, \(\Pi'\subset H\), and
\(\Pi''\subset[H,H]\subset N\).  Moreover
\(Q_\infty\) is an inverse limit of finite two-groups
with abelian normal module and cyclic quotient.  Its
universal map from \(\Pi\) therefore factors through
\(\Pi_2/\Pi_2''\).  Both kernels used in (A.157.5)
are characteristic in the respective geometric group,
so these are maps of the actual arithmetic extensions.

\emph{A fixed sheet gives the precise arithmetic lift.}
Let \(E_D=\pi_1(D^\circ)\subset E_U\), using compatible
geometric base points.  No rational point of \(D^\circ\)
is needed for this inclusion.  The finite scheme obtained
by pulling the \(\mu_8\)-torsor (A.157.1) back along a
section of \(E_U/N\) has cohomology class
\(2\kappa_8(t)+\kappa_8(1-t)-[64a^2]\).
The character descends to this quotient because
\(N\subset H\).  Equation (A.157.4) says exactly that
this torsor is trivial, and hence it has a rational
point, or equivalently a fixed geometric sheet.
The geometric group acts transitively on the sheets.
After a geometric conjugation the given section
therefore takes values in
\[
                         E_D/N\subset E_U/N.
\tag{A.157.14}
\]

Over \(K\), compactification and the literal positive
Prym morphism of A.150 give a homomorphism
\[
 (E_D/N)\times_{G_\Q}G_K
                      \longrightarrow T_2E_b\rtimes G_K.
\tag{A.157.15}
\]
To verify its factorization, the compactification
first kills \(\mathcal I\); the geometric part of
the elliptic two-primary quotient is abelian and
pro-two, so the proper map kills the remaining
kernel in (A.157.2).  The splitting on the target
comes from its actual origin.  A chosen geometric
path between the image base point and that origin
fixes the homomorphism; changing the path changes
the resulting section cocycle by a coboundary.
The conjugate Prym map over \(K\) also kills the
same \(N\), by the same proper abelian quotient.
All arithmetic normality was established before
choosing either Prym map.

\emph{Every local lift is an actual point or cusp lift.}
Fix a place \(p\) of \(\Q\), including infinity.
By the assumed local matching, choose a full ordinary
or cuspidal local section \(s_p\) of \(E_U\), whose
image is conjugate to the localization in (A.157.14).
Lift its geometric conjugator from \(\Pi/N\) to
\(\Pi\).  The conjugated full local section has image
in \(E_D\): the inverse image of \(E_D/N\) in \(E_U\)
is exactly \(E_D\), since \(N\subset H\subset E_D\).
Its projection is the localization of (A.157.14).

For an ordinary section, the fixed sheet is an actual
\(\Q_p\)-point of the finite etale fibre above its
centre.  This gives a point of \(D^\circ(\Q_p)\).
For a cuspidal section, its fixed sheet belongs to a
boundary inertia orbit.  That orbit is stable under
\(G_{\Q_p}\), so the corresponding point of the
proper normalization is rational over \(\Q_p\).
The intersection of the original cusp decomposition
group with \(E_D\) is the decomposition group of
that boundary point.  Compactification kills its
boundary inertia, and its proper section is the
section of this actual point.  The same argument
at infinity identifies a fixed ordinary sheet with
a real point, and a stable inertia orbit with a
real boundary point.  Thus every local map in
(A.157.15) is the actual Kummer class of a proper
elliptic point.  Base change to each completion
\(K_w\) preserves these point and decomposition-group
maps.

The section obtained from (A.157.15) consequently
has finite classes in every
\(\operatorname{Sel}^{2^n}(E_b/K)\), compatibly in
\(n\).  A.154.28--33 applies to exactly these finite
classes and their continuous cocycle: its image is
the Kummer class of one global point
\(Q\in\{O,(0,0)\}\).  At two the ordinary local
matching supplies an actual point
\(P\in D^\circ(\Q_2)\).  The injective dyadic
Kummer map and the exact fibres in A.152 force
\(P\) to be one of the four deleted boundary
points.  This is the required contradiction.
It proves the quotient assertion and, through the
literal maps (A.157.5), both metabelian assertions.

\emph{Completed coordinates and the proper rational points.}
For a locally matching full metabelian section with
ordinary completed coordinates \((A,B)\), the local
full Kummer coordinates are \((x_p,1-x_p)\).
At two, a cusp of label zero would give \(B_2=1\),
a cusp of label one would give \(A_2=1\), and a
cusp at infinity would give \(B_2/A_2=-1\).
The injective local multiplicative Kummer map turns
these into \(x_2=0\), \(x_2=1\), or \(1=0\),
respectively.  All contradict ordinariness.
An ordinary local representative has the same centre
by that injection.  The assumed global combined
class is (A.157.4), so the proved exclusion applies.
In the pro-two metabelian quotient the same argument
at the physical place two uses its own-coefficient
Kummer injection from A.97: \(\Q_2\) has no
nontrivial roots of unity of odd order.

Finally every \(P\in\overline D(K)\) maps under
\(q_+\) to \(E_b(K)=\{O,(0,0)\}\) by A.154.
The exact degree-two fibres in A.152 are precisely
the four displayed boundary points, all defined
over \(\Q\).  This proves (A.157.6).
\end{proof}

For the finite comparisons used in A.158, set
\[
 N_m=\ker\{H\longrightarrow(H/\mathcal I)^{\rm ab}/2^m\},
 \qquad
 Q_m=\Pi/N_m\simeq
 \bigl((\Z/2^m)[V]/(V^6+V^4+V^2+1)\bigr)\rtimes C_8.
\tag{A.157.16}
\]
The same arithmetic normality and the same inverse maps
prove this equality.  Its order is \(2^{6m+3}\),
and \(N=\bigcap_mN_m\).  At \(m=1\), substitution
\(V=1+\tau\) gives \(V^6+V^4+V^2+1=\tau^6\),
so this is the literal order-512 quotient of A.141.
After restriction to \(K\), (A.157.15) projects to
\[
       (E_D/N_m)\times_{G_\Q}G_K
                       \longrightarrow E_b[2^m]\rtimes G_K.
\tag{A.157.17}
\]
These are the actual finite multiplication-quotient
maps used for the point-dependent finite test.
The exact integral inertia vectors, their primitive
relations, the companion action, the finite reductions,
and the two curve-map compositions are checked in
\path{code/harmonic_metabelian_class_certificate.py}.
